% =====================================================================
%  Apêndice F — Registro da campanha de coleta (novo na AP2)
% =====================================================================
\chapter{Registro da campanha de coleta}
\label{ap:campanha}

\section*{Diário da campanha}
\phantomsection
\addcontentsline{toc}{section}{Diário da campanha}

\begin{table}[H]
\centering
\caption{Diário da campanha de coleta}
\label{tab:diario}
\footnotesize
\setlength{\tabcolsep}{3.5pt}
\begin{tabular}{@{}L{2.4cm}L{4.6cm}L{3.4cm}L{3.6cm}@{}}
\toprule
\textbf{Data e hora} & \textbf{O que aconteceu} & \textbf{Como a equipe soube} & \textbf{O que a equipe fez} \\
\midrule
24/10 14:07 & nó instalado num quarto, junto a uma janela aberta, a \SI{50}{\centi\metre} do chão; gravador iniciado no notebook, a cerca de \SI{5}{\metre} & --- & --- \\
24/10 14:17 & \CT{15}: Bluetooth do notebook desligado por \SI{2}{\minute} & diário automático e \texttt{eventos.csv} & nada: o gravador reconectou sozinho às 14:19:49 \\
24/10 14:34 & \CT{16} automático falhou: o \texttt{kill -INT} foi ignorado e o gravador seguiu gravando & conferência às 15:10, nas leituras que continuaram chegando & corrigiu o \texttt{gravador.py} e trocou o processo às 15:11 \\
24/10 15:11 & na troca, o BlueZ manteve a conexão do gravador antigo; o nó parou de anunciar & \texttt{nao\_encontrado} no terminal & \texttt{bluetoothctl disconnect}; reconectou às 15:12:44, sem perda \\
24/10 15:18 a 16:11 & \CT{18}: USB do nó arrancado 4 vezes (15:18, 15:20, 15:22, 16:11) & hora anotada no terminal; \texttt{eventos.csv} & religou o USB; o gravador reconectou sozinho nas quatro \\
24/10 15:25 & \CT{17}: botão EN & hora anotada no terminal & nada: reconexão em \SI{5.6}{\second} \\
24/10 15:31 a 16:06 & \CT{16} refeito: gravador parado por \SI{35}{\minute} & diário automático & gravador de volta às 16:06:18 \\
25/10 00:28 & nó reposicionado (mesma altura, junto à janela, a cerca de \SI{5}{\metre}) & anotação da equipe & nada: a conexão não caiu \\
25/10 09:05 & nó de volta à posição original & anotação da equipe & --- \\
25/10 14:12 & fim da coleta (gravador encerrado às 14:12:37) & --- & \texttt{qualidade.py} sobre o arquivo final \\
\bottomrule
\end{tabular}
\fonte{Elaborada pela equipe.}
\end{table}

\section*{Registro de eventos para rotulagem}
\phantomsection
\addcontentsline{toc}{section}{Registro de eventos para rotulagem}

O que aconteceu em volta do nó enquanto ele media, com hora: o sol que
entrou direto no sensor, a sombra, a janela aberta ou fechada, a chuva,
alguém mexendo no nó. É deste registro que sai a
rotulagem no Módulo 3.

\begin{table}[H]
\centering
\caption{Eventos observados durante a coleta}
\label{tab:eventos}
\footnotesize
\begin{tabular}{@{}L{2.6cm}L{2.6cm}L{5.4cm}L{3.4cm}@{}}
\toprule
\textbf{Início} & \textbf{Fim} & \textbf{Evento} & \textbf{Quem anotou, e como} \\
\midrule
24/10 14:07 & 24/10 14:55 & sol direto no sensor, janela aberta; DHT11 aquecido pelo sol (até \SI{44.4}{\celsius}) & equipe, conferido nos dados \\
24/10 14:55 & 24/10 17:53 & sem sol direto; luz natural caindo até o anoitecer & equipe, conferido nos dados \\
24/10 17:53 & 25/10 00:28 & escuro lá fora; lâmpada do quarto acesa (\SI{15}{lux}) & equipe \\
24/10 20:34 & 25/10 09:05 & janela e cortinas fechadas; ar-condicionado ligado & equipe \\
25/10 00:28 & 25/10 09:05 & luz apagada; nó reposicionado no quarto & equipe \\
25/10 09:05 & 25/10 14:13 & cortinas e janela abertas, ar-condicionado desligado, nó na posição original; sol direto a partir de cerca de 13:55 & equipe, conferido nos dados \\
\bottomrule
\end{tabular}
\fonte{Elaborada pela equipe.}
\end{table}

\section*{Saída do relatório de qualidade}
\phantomsection
\addcontentsline{toc}{section}{Saída do relatório de qualidade}

A saída completa do \texttt{qualidade.py} sobre o arquivo final, sem
edição, gerada pelo comando abaixo, rodado na raiz do repositório. O
quadro lê \caminho{docs/figuras/qualidade.txt}; enquanto o arquivo não
existe, mostra um aviso de pendência.

\begin{lstlisting}[style=terminal, numbers=none]
python gravador/qualidade.py dados/mensagens.csv --eventos dados/eventos.csv \
    --inicio '2026-10-24 14:07' --fim '2026-10-25 14:13' > dados/qualidade.txt
cp dados/qualidade.txt docs/figuras/qualidade.txt
\end{lstlisting}

\IfFileExists{figuras/qualidade.txt}%
  {\def\arquivoQualidade{figuras/qualidade.txt}}%
  {\def\arquivoQualidade{figuras/qualidade-pendente.txt}}
\lstinputlisting[style=terminal,
  caption={Saída do \texttt{qualidade.py} sobre o arquivo final da campanha},
  label={lst:qualidade}]{\arquivoQualidade}
